\documentclass[12pt]{article}
\usepackage[a4paper,margin=1in]{geometry}
\usepackage[utf8]{inputenc}
\usepackage[T1]{fontenc}
\usepackage{amsmath,amssymb}
\usepackage{textcomp}
\usepackage{parskip}
\begin{document}
\begin{center}
{\Large \textbf{2002 STEP II}}
\end{center}

\subsection*{Page 1}
\begin{flushleft}
STEP II, 2002 2\par
\bigskip\textbf{Section A: Pure Mathematics}\par
\bigskip\textbf{Question 1.} Show that\par
\[\pi/4\]
\[1 \sqrt\{3\} 1\]
\[dθ = .\]
1 cos 2θ 2 - 2\par
\[Z\pi/6\]
-\par
\[By using the substitution x = sin 2θ, or otherwise, show that\]
\[1 \pi\]
\[dx = \sqrt\{3\} 1 .\]
\[1 \sqrt(1 x^2) - - 6\]
\[Z\sqrt\{3\}/2\]
- -\par
Hence evaluate the integral\par
\[2/\sqrt\{3\}\]
dy .\par
\[y(y \sqrt(y^2 12))\]
Z - -\par
\bigskip\textbf{Question 2.} Show that setting z z 1 = w in the quartic equation\par
-\par
-\par
4 3 2\par
\[z + 5z + 4z 5z + 1 = 0\]
-\par
\[results in the quadratic equation w^2 + 5w + 6 = 0. Hence solve the above quartic equation.\]
Solve similarly the equation\par
8 7 6 5 4 3 2\par
\[2z 3z 12z + 12z + 22z 12z 12z + 3z + 2 = 0 .\]
- - - -\par
\bigskip\textbf{Question 3.} The nth Fermat number, F , is defined by\par
n\par
2n\par
\[F = 2 + 1 , n = 0, 1, 2, . . . ,\]
n\par
where 2\par
2n\par
\[means 2 raised to the power 2n . Calculate F^0, F^1, F^2 and F^3 . Show that, for k = 1,\]
\[k = 2 and k = 3 ,\]
\[F0F1 . . . F k 1 = F k 2 . ( )\]
- - ∗\par
\[Prove, by induction, or otherwise, that ( ) holds for all k > 1. Deduce that no two Fermat\]
∗\par
numbers have a common factor greater than 1.\par
Hence show that there are infinitely many prime numbers.\par
\end{flushleft}
\subsection*{Page 2}
\begin{flushleft}
STEP II, 2002 3\par
q\par
\[4 Give a sketch to show that, if f (x) > 0 for p < x < q , then f (x) dx > 0 .\]
p\par
R\par
\textbf{(i)} By considering f (x) = ax2 bx+c show that, if a > 0 and b2 < 4ac, then 3b < 2a+6c .\par
-\par
\textbf{(ii)} By considering f (x) = a sin 2 x b sin x + c show that, if a > 0 and b2 < 4ac, then\par
-\par
\[4b < (a + 2c)\pi .\]
\textbf{(iii)} Show that, if a > 0, b2 < 4ac and q > p > 0 , then\par
1 1\par
\[b ln(q/p) < a + c(q p) .\]
p - q -\par
\lambda \mu\par
\bigskip\textbf{Question 5.} The numbers x , where n = 0, 1, 2, . . . , satisfy\par
n\par
\[x n+1 = kx n (1 x n ) .\]
-\par
\textbf{(i)} Prove that, if 0 < k < 4 and 0 < x0 < 1, then 0 < x\par
n\par
\[< 1 for all n .\]
\textbf{(ii)} Given that x0 = x1 = x2 = = a , with a = 0 and a = 1, find k in terms of a .\par
\cdots 6 6\par
\textbf{(iii)} Given instead that x0 = x2 = x4 = = a , with a = 0 and a = 1, show that\par
ab3 b2 + (1 a) = 0, where b = k(1 · · a · ) . Given, in add 6 ition, that 6 x1 = a, find the\par
- - - 6\par
possible values of k in terms of a .\par
\bigskip\textbf{Question 6.} The lines l1, l2 and l3 lie in an inclined plane P and pass through a common point A. The line\par
l2 is a line of greatest slope in P . The line l1 is perpendicular to l3 and makes an acute angle α\par
with l2. The angles between the horizontal and l1, l2 and l3 are \pi/6, β and \pi/4, respectively.\par
\[Show that cos α sin β = and find the value of sin α sin β . Deduce that β = \pi/3 .\]
The lines l1 and l3 are rotated in P about A so that l1 and l3 remain perpendicular to each\par
other. The new acute angle between l1 and l2 is θ. The new angles which l1 and l3 make with\par
the horizontal are φ and 2φ, respectively. Show that\par
\[3 + \sqrt\{13\}\]
\[tan θ = .\]
\end{flushleft}
\subsection*{Page 3}
\begin{flushleft}
STEP II, 2002 4\par
\bigskip\textbf{Question 7.} In 3-dimensional space, the lines m1 and m2 pass through the origin and have directions i + j\par
and i + k, respectively. Find the directions of the two lines m3 and m4 that pass through the\par
origin and make angles of \pi/4 with both m1 and m2. Find also the cosine of the acute angle\par
\[between m^3 and m^4.\]
The points A and B lie on m1 and m2 respectively, and are each at distance λ\sqrt2 units from O.\par
The points P and Q lie on m3 and m4 respectively, and are each at distance 1 unit from O. If\par
all the coordinates (with respect to axes i, j and k) of A, B, P and Q are non-negative, prove\par
that:\par
\textbf{(i)} there are only two values of λ for which AQ is perpendicular to BP ;\par
\textbf{(ii)} there are no non-zero values of λ for which AQ and BP intersect.\par
\bigskip\textbf{Question 8.} Find y in terms of x, given that:\par
dy\par
\[for x < 0 , = y and y = a when x = 1 ;\]
dx - -\par
dy\par
\[for x > 0 , = y and y = b when x = 1 .\]
dx\par
Sketch a solution curve. Determine the condition on a and b for the solution curve to be\par
\[continuous (that is, for there to be no 'jump' in the value of y) at x = 0.\]
Solve the differential equation\par
dy\par
\[= ex 1 y\]
dx | - |\par
\[given that y = e e when x = 1 and that y is continuous at x = 0 . Write down the following\]
limits:\par
\textbf{(i)} lim y exp( ex) ; (ii) lim y e x .\par
-\par
x + - x\par
\to \infty \to-\infty\par
\end{flushleft}
\subsection*{Page 4}
\begin{flushleft}
STEP II, 2002 5\par
\bigskip\textbf{Section B: Mechanics}\par
\bigskip\textbf{Question 9.} A particle is projected from a point O on a horizontal plane with speed V and at an angle\par
of elevation α. The vertical plane in which the motion takes place is perpendicular to two\par
vertical walls, both of height h, at distances a and b from O. Given that the particle just\par
passes over the walls, find tan α in terms of a, b and h and show that\par
2V 2 ab (a + b)2h\par
\[= + .\]
g h ab\par
The heights of the walls are now increased by the same small positive amount δh . A second\par
particle is projected so that it just passes over both walls, and the new angle and speed of\par
projection are α + δα and V + δV , respectively. Show that\par
a + b\par
sec α δα δh ,\par
\approx ab\par
\[and deduce that δα > 0 . Show also that δV is positive if h > ab/(a + b) and negative if\]
\[h < ab/(a + b) .\]
\bigskip\textbf{Question 10.} A competitor in a Marathon of 42 3 km runs the first t hours of the race at a constant speed\par
of 13 km h 1 and the remainder at a constant speed of 14 + 2t/T km h 1, where T hours is\par
- -\par
her time for the race. Show that the minimum possible value of T over all possible values of\par
t is 3.\par
The speed of another competitor decreases linearly with respect to time from 16 km h 1 at the\par
-\par
start of the race. If both of these competitors have a run time of 3 hours, find the maximum\par
distance between them at any stage of the race.\par
\bigskip\textbf{Question 11.} A rigid straight beam AB has length l and weight W . Its weight per unit length at a distance\par
x from B is αW l 1(x/l)α 1 , where α is a positive constant. Show that the centre of mass of\par
- -\par
the beam is at a distance αl/(α + 1) from B.\par
The beam is placed with the end A on a rough horizontal floor and the end B resting against\par
a rough vertical wall. The beam is in a vertical plane at right angles to the plane of the wall\par
and makes an angle of θ with the floor. The coefficient of friction between the floor and the\par
beam is µ and the coefficient of friction between the wall and the beam is also µ . Show that,\par
if the equilibrium is limiting at both A and B, then\par
1 αµ2\par
\[tan θ = - .\]
(1 + α)µ\par
\[Given that α = 3/2 and given also that the beam slides for any θ < \pi/4 find the greatest\]
possible value of µ .\par
\end{flushleft}
\subsection*{Page 5}
\begin{flushleft}
STEP II, 2002 6\par
\bigskip\textbf{Section C: Probability and Statistics}\par
\bigskip\textbf{Question 12.} On K consecutive days each of L identical coins is thrown M times. For each coin, the\par
\[probability of throwing a head in any one throw is p (where 0 < p < 1). Show that the\]
probability that on exactly k of these days more than l of the coins will each produce fewer\par
than m heads can be approximated by\par
K\par
qk(1 q)K k,\par
-\par
k -\par
\lambda \mu\par
where\par
2h 2l 1 2m 1 2M p\par
\[q = Φ - - , h = LΦ - -\]
\[2\sqrt\{h\} 2 M p(1 p) !\]
\lambda \mu -\par
and Φ(.) is the cumulative distribution function of a standpard normal variate.\par
Would you expect this approximation to be accurate in the case K = 7, k = 2, L = 500, l = 4,\par
\[M = 100, m = 48 and p = 0.6 ?\]
\bigskip\textbf{Question 13.} Let F(x) be the cumulative distribution function of a random variable X, which satisfies\par
\[F(a) = 0 and F(b) = 1, where a > 0. Let\]
F(y)\par
\[G(y) = .\]
2 F(y)\par
-\par
\[Show that G(a) = 0 , G(b) = 1 and that G (y) > 0 . Show also that\]
1 2\par
6 6 2 .\par
2 (2 F(y))2\par
-\par
The random variable Y has cumulative distribution function G(y) . Show that\par
1 E(X) 6 E(Y ) 6 2E(X) ,\par
and that\par
Var (Y ) 6 2 Var (X) + 7 (E(X)) 2 .\par
\end{flushleft}
\subsection*{Page 6}
\begin{flushleft}
STEP II, 2002 7\par
\bigskip\textbf{Question 14.} A densely populated circular island is divided into N concentric regions R1, R2, . . . , R\par
N\par
, such\par
that the inner and outer radii of R are n 1 km and n km, respectively. The average number\par
n\par
-\par
of road accidents that occur in any one day in R is 2 n/N , independently of the number\par
n\par
-\par
of accidents in any other region.\par
Each day an observer selects a region at random, with a probability that is proportional to\par
the area of the region, and records the number of road accidents, X, that occur in it. Show\par
that, in the long term, the average number of recorded accidents per day will be\par
1 1 1\par
2 1 + 4 .\par
- 6 N - N\par
\lambda \mu \lambda \mu\par
N\par
\[[Note: n^2 = 1 N (N + 1)(2N + 1) .]\]
\[n=1\]
P\par
Show also that\par
e 2N k 2 N\par
\[P(X = k) = - - - (2n 1)(2N n)k en/N .\]
k! - -\par
\[n=1\]
X\par
\[Suppose now that N = 3 and that, on a particular day, two accidents were recorded. Show\]
\[that the probability that R^2 had been selected is\]
.\par
\[48 + 45 e^1/3 + 25 e 1/3\]
-\par
\end{flushleft}
\end{document}
